\begin{abstract}
\textbf{Objective:} Covariance estimation, predictive
Gaussian likelihood, and held-source-participant
discrimination may select different EEG
regularisation regimes. We examine whether
these differences affect zero-calibration
cross-participant decoding, and when a
parameter change becomes an actual
classification-ranking change.
\textbf{Methods:} We retain the original
source-fitted covariance/tangent-LDA pathway,
subject-grouped validation, and frozen
historical tests. We distinguish validation
class-centred covariance scoring from genuine
source-mean predictive likelihood, then
decompose AUC changes exactly into
positive--negative pair transitions. A spectral
interval enclosure and second-order Taylor
refinement bound rank stability for every
shrinkage value within a continuous interval.
A source-only branch-and-bound search returns
an explicit finite-validation optimisation
gap or reports failure.
\textbf{Results:} In the prospective EEGMMIDB
20-person test block, source-discriminative
selection exceeded the historical class-centred
selector by $0.0561$ AUC (Holm $p=0.00597$),
but harmed five participants by $>0.02$ AUC.
Corrected predictive-likelihood comparisons
narrowed the apparent advantage in other
cohorts. The new 3-channel BNCI2014\_004
CVaR control altered 8/9 shrinkage decisions
without improving mean held-person AUC.
On BNCI2014\_002, a retrospective 14-pool
\emph{source-validation} audit showed
mean AUC difference $+0.02680$ and
spectrally certified continuous-path
pair-rank stability of $32.48\%$ after
Taylor refinement. A five-split continuous
source search failed its declared $0.02$
optimality tolerance (residual gap $0.08643$).
\textbf{Conclusion:} Regularisation-objective
disagreement need not imply transfer gain.
Continuous spectral analysis makes the
source-domain ranking consequences auditable,
but neither its bound nor the new search
establishes improved unknown-participant
performance. Clinical generalisation
requires a genuinely external frozen test.
\end{abstract}
